\paragraph{Definitions and notations (original Sections 1.1, 2.1).}
For a complexity class $\mathfrak C$, a complexity-theoretic statement is $\mathfrak C$-relativizing if it is true relative to every oracle $\mathcal O\in\mathfrak C$.
\paragraph{Definition 20.} A language $L$ is in $\mathsf{EXPH}$ if there is a constant $k$, an exponential bound $e(n)=2^{\operatorname{poly}(n)}$, and a $\operatorname{poly}(e(n))$-time algorithm $M$ such that, for every input $x\in\{0,1\}^*$,
\[x\in L\Longleftrightarrow\exists y_1\forall y_2\exists y_3\forall y_4\ldots\exists y_k,\ M(x,y_1,\ldots,y_k)=1,\]
where $y_i\in\{0,1\}^{e(n)}$ for each $i\in[k]$. Equivalently, a language is in $\mathsf{EXPH}$ if it can be decided in exponential ($2^{\operatorname{poly}(n)}$) time with access to a $\mathsf{PH}$ oracle.
\paragraph{Pseudodeterministic algorithms.} A pseudodeterministic algorithm is a probabilistic algorithm that returns a fixed output with high probability on any given input. We will also consider pseudodeterministic oracle (or advice-taking) algorithms. Such an algorithm will maintain its pseudodeterministic behavior (i.e., outputting a fixed answer with high probability for any input) when a correct oracle (or advice) is given.
\paragraph{Dense property (original Section 1.3.2).} A property $Q\subseteq\{0,1\}^*$ is dense if for every input length $n\in\mathbb N$, we have that $|Q\cap\{0,1\}^n|\geq2^n/\operatorname{poly}(n)$.
\paragraph{Theorem 23 [68].} For every oracle $\mathcal O$, $\mathsf{BPP}^{\mathcal O}\subseteq(\Sigma^p_2)^{\mathcal O}$.
